% =====================================================================
\lezione{26/10/2026}{Risk management ISO e introduzione al Cyber Resilience Act}
  {Gianluca Gilardi; Prof.~Giuseppe Vaciago}
  {minacce e threat intelligence; scale di probabilità e impatto; prioritizzazione; metodologia ENISA
   per i dati personali; ISO vs ENISA; risk register; opzioni di trattamento; Statement of
   Applicability; piano di trattamento; casi di digital forensics su dispositivi smart; CRA: ambito,
   esclusioni, requisiti essenziali, prodotti importanti e critici, obblighi della filiera, segnalazione;
   briefing del laboratorio~3}
\label{lez:2026-10-26}
% =====================================================================

\begin{nota}[Fonte]
Trascrizione automatica della registrazione, revisionata (inizio con prove audio; pausa intorno ai
minuti 63--78). Slide di riferimento: \emph{06.~27001\_2026\_2310} (\emph{Risk Management}) per la prima
parte, \emph{05.~PoliTO\_CRA} per la seconda; briefing del laboratorio in
\texttt{lab/lab\_03/materiale}. La parte del CRA prosegue e si conclude nella lezione del 02/11
(Capitolo~\ref{lez:2026-11-02}).
\end{nota}

\section{Sintesi}
Tre blocchi. Gianluca completa il \textbf{risk management secondo la ISO}: scale, prioritizzazione,
confronto con la metodologia ENISA per i dati personali, contenuto del risk register, opzioni di
trattamento, Statement of Applicability e piano di trattamento. Il professore introduce poi il
\textbf{Cyber Resilience Act} partendo da casi in cui dispositivi smart sono diventati prove in giudizio,
per mostrare quanto IoT e dispositivi indossabili siano sottovalutati sul piano della sicurezza; segue la
prima metà delle slide (ambito, esclusioni, requisiti, filiera, segnalazione). Chiude il briefing del
laboratorio~3 sul risk register.

% ---------------------------------------------------------------------
\section{Organizzazione}
\begin{itemize}
  \item \textbf{Laboratorio~3, venerdì 30/10}: compilazione di un risk register ISO~27001. Materiale nella
    cartella della piattaforma ``lab 2026 1030'': briefing (PDF), playbook (Markdown) e una copia di prova
    del file Excel senza scenari, per esercitarsi. Il playbook va letto prima.
  \item Domande sull'esercitazione: per email a Gianluca, o in aula il venerdì.
  \item Il professore resta disponibile durante la pausa per chi vuole discutere il laboratorio~2.
  \item La parte finale del CRA (slide residue e un ``meme'' sulla filiera) è rinviata al lunedì
    successivo.
\end{itemize}

% =====================================================================
\section{Risk management ISO (Gianluca)}
% =====================================================================

\subsection{Minacce, vulnerabilità e threat intelligence}
Le minacce sfruttano vulnerabilità tecniche (individuabili con penetration test e scansioni) o umane
(ingegneria sociale). L'ultima revisione della 27001 chiede di individuare le proprie \textbf{fonti di
threat intelligence}: non si deve conoscere ogni minaccia da soli, ma attingere a fonti esterne, come i
bollettini dell'ACN. La probabilità di sfruttamento di molte vulnerabilità note è già stimata nelle fonti
di threat intelligence.

\subsection{Scale, prioritizzazione, propensione al rischio}
\begin{itemize}
  \item Le matrici sono quasi sempre \textbf{qualitative}: alto, medio, basso (al massimo quattro livelli,
    con ``molto alto''); le scale quantitative sono rare.
  \item Formula di base: \textbf{Rischio = Impatto $\times$ Probabilità}.
  \item Prioritizzare: di solito si parte dal rischio complessivo, ma conta anche il tempo di attuazione
    delle contromisure. Le \emph{quick win} possono passare avanti a rischi più alti che richiedono
    progetti lunghi. Decide il vertice (il board), che però chiede a chi valuta costi, tempi e proposta di
    priorità.
  \item La ISO \textbf{non fissa} alcuna matrice: si deve essere coerenti e saper difendere le scelte. La
    più usata è la \textbf{3$\times$3}.
  \item \textbf{Risk appetite} (slide): disponibilità complessiva ad assumere rischi, decisione strategica
    del board. \textbf{Risk tolerance}: livello accettato per singole attività o processi, livello
    operativo con soglie specifiche.
\end{itemize}

\paragraph{Livelli di impatto sull'organizzazione (indicazioni del docente).} Basso: perdita economica
limitata, nessun fermo operativo, recupero facile. Medio: un ``mal di testa'' misurabile (sanzioni, fermo
limitato nel tempo) ma superabile. Alto (o molto alto): evento che può mettere a rischio
\textbf{l'esistenza stessa} dell'organizzazione, come la perdita totale dei dati senza backup.

\subsection{La metodologia ENISA per i dati personali}
ENISA non ha pubblicato solo la metodologia di severità dei breach, ma anche linee guida per le PMI sulla
valutazione del rischio dei trattamenti. Sono pensate per i dati personali, ma ogni dato personale è
un'informazione: il metodo è un buon punto di partenza per costruirne uno proprio.
\begin{enumerate}
  \item \textbf{Contesto del trattamento}: quali dati, finalità, mezzi, luogo (on premise o cloud, con
    rischi di terze parti), categorie di interessati, destinatari; tutte le fasi: raccolta, conservazione,
    uso, trasferimento, cancellazione. Output: il perimetro della valutazione.
  \item \textbf{Impatto sugli individui}, su quattro livelli (basso, medio, alto, molto alto), valutato
    separatamente per perdita di riservatezza, integrità e disponibilità: \textbf{si prende il più alto}.
  \item \textbf{Probabilità della minaccia}, su quattro aree (risorse di rete e tecniche; processi e
    procedure; persone e terze parti; settore e dimensione), ciascuna con punteggio 1--3. Totale: 4--5
    bassa, 6--8 media, 9--12 alta. Le definizioni sono qualitative (``improbabile'', ``ragionevole
    possibilità'', ``probabile''), senza percentuali.
  \item \textbf{Valutazione del rischio} con la matrice probabilità $\times$ impatto.
\end{enumerate}

\begin{center}
\small
\textbf{Matrice ENISA (slide)}\par\smallskip
\begin{tabular}{llll}
\toprule
\textbf{Probabilità} & \textbf{Impatto basso} & \textbf{Impatto medio} & \textbf{Impatto alto/molto alto} \\ \midrule
Bassa & basso & medio & alto \\
Media & medio & alto & alto \\
Alta  & alto & alto & alto \\ \bottomrule
\end{tabular}
\end{center}

\subsection{ISO 27001 e ENISA a confronto}
\begin{center}
\small
\begin{tabularx}{\linewidth}{l X X}
\toprule
\textbf{Aspetto} & \textbf{ISO/IEC 27001} & \textbf{ENISA (dati personali)} \\ \midrule
Ambito & sicurezza delle informazioni in generale & trattamenti di dati personali \\
Impatto & sull'organizzazione (triade CIA) & sui diritti e le libertà degli individui \\
Calcolo & metodologie libere & approccio strutturato in 4 passi \\
Minacce & categorie generali & questionari per contesto \\
Controlli & 93 controlli (Allegato A) & misure proporzionate al livello di rischio \\
Conformità & certificazione ISO & art.~32 GDPR \\ \bottomrule
\end{tabularx}
\end{center}
I due approcci si usano spesso insieme: ENISA per i trattamenti di dati personali, ISO per l'ISMS, un
unico registro con doppia valutazione. La metodologia ENISA è apprezzata anche nella comunità ISO.

\subsection{Il risk register}
Non esiste un formato obbligatorio (Excel o piattaforme SaaS). Campi che, per esperienza, vanno considerati
irrinunciabili:
\begin{itemize}
  \item \textbf{ID univoco} del rischio, per tracciarlo in tutti i documenti e progetti;
  \item asset coinvolti, minaccia, vulnerabilità;
  \item probabilità e impatto \textbf{inerenti} e, dopo i controlli, \textbf{residui};
  \item livello di rischio calcolato con la matrice;
  \item \textbf{risk owner}: chi decide e firma il trattamento e ne risponde;
  \item trattamento proposto (misure già adottate e da adottare).
\end{itemize}

\subsection{Opzioni di trattamento e controlli}
\begin{itemize}
  \item \textbf{Mitigare} (ridurre): controlli che riducono probabilità o impatto; è l'opzione più comune.
  \item \textbf{Accettare}: il rischio rientra nella propensione al rischio (vale per l'inerente o per il
    residuo); richiede approvazione formale.
  \item \textbf{Evitare}: rinunciare all'attività che genera il rischio.
  \item \textbf{Trasferire}: condividere il rischio con un terzo (assicurazione, outsourcing). Non elimina
    la responsabilità: anche il trasferimento è una decisione di cui si risponde e che deve essere coerente
    con la propensione al rischio. La \emph{shared responsibility matrix} dei cloud provider è una forma di
    trasferimento parziale.
\end{itemize}
La 27001 dà i nomi dei controlli, la 27002 i mezzi per attuarli (es.\ controllo d'accesso attuato con
l'MFA). \textbf{Tipi di controllo}: preventivi (VA, penetration test), detective (IDS/IPS, DLP),
correttivi (disaster recovery, business continuity).

\subsection{Statement of Applicability e piano di trattamento}
\paragraph{SoA.} Documento \textbf{obbligatorio} per la conformità 27001: elenca tutti i controlli
dell'Allegato~A, dichiara se ciascuno è applicabile, ne giustifica inclusione o esclusione e ne indica lo
stato di attuazione. Documenta le decisioni prese nel trattamento del rischio. Per Gianluca è ancora più
importante saper giustificare perché un controllo \textbf{non} è attuato; l'auditor può contestare il
legame tra un controllo e un rischio e imporre di correggere la SoA.

\paragraph{Piano di trattamento.} Per ogni rischio da trattare: opzione scelta, azioni concrete, owner di
ogni azione (accountability), tempi di completamento, risorse (budget e personale: il management chiederà
sempre ``quanto costa?''), giustificazione. La giustificazione discende in gran parte dalla propensione al
rischio: si accetta perché sotto soglia; si mitiga perché l'inerente è sopra soglia e i controlli portano il
residuo entro la soglia; si trasferisce ciò che resta fuori dal proprio controllo; si evita quando nessun
residuo è accettabile. La matrice \textbf{RACI} (responsabile, accountable, consultato, informato) è
fondamentale nella procedura di gestione dei data breach, ma non è richiesta per il piano di trattamento.

\paragraph{Monitoraggio continuo.} I rischi cambiano, nascono e muoiono: il registro non si chiude in un
cassetto. La threat intelligence dimostra all'auditor l'attenzione continua al rischio.

\begin{esame}
\textbf{Messaggi chiave del risk management (slide finali):} è il fondamento della cybersicurezza; servono
asset, minacce, vulnerabilità, impatto e probabilità; si segue un processo strutturato (ISO, NIST, FAIR,
ENISA: non serve reinventare la ruota); il trattamento è \textbf{strategico}, non solo tattico (non ha
senso spendere per mitigare un processo che verrà dismesso a breve); le decisioni si documentano nella
\textbf{SoA}; il monitoraggio è continuo.
\end{esame}

% =====================================================================
\section{Cyber Resilience Act (Prof.~Vaciago)}
% =====================================================================

\subsection{Introduzione: i dispositivi smart come fonte di prova e di rischio}
Fuori programma, per introdurre il tema. Ogni dispositivo digitale tocca due ambiti: \textbf{digital
forensics} e \textbf{cybersecurity}. Le stesse informazioni utili alla polizia sono utili a chi vuole
commettere un reato.
\begin{itemize}
  \item \textbf{Caso Fitbit}: in un'indagine per omicidio i dati dei braccialetti dei coniugi hanno
    smentito la versione del marito, che ha poi confessato.
  \item \textbf{Caso Amazon Echo}: la polizia cercava registrazioni ambientali; Amazon ha chiarito che non
    registra in continuo (solo dopo la parola di attivazione o per attivazioni accidentali). L'aspetto
    curioso: inizialmente ha negato i dati sostenendo che solo il proprietario potesse chiederli.
  \item \textbf{Smart TV}: studi recenti (il professore cita la ricercatrice Anna Maria Mandalari)
    mostrano funzioni nascoste che attivano la telecamera per rilevare le reazioni degli spettatori;
    erano previste nelle condizioni d'uso.
  \item \textbf{Caso Ross Compton}: accusato di aver incendiato la propria casa per frodare
    l'assicurazione, è stato condannato anche grazie ai dati del suo pacemaker (battito incompatibile con
    la sua versione). Dibattito sull'uso processuale dei dati sanitari.
\end{itemize}
Il professore richiama il rischio legato alla quantità di dati sanitari ceduti spontaneamente a
dispositivi e assistenti di IA, spesso in mano a soggetti statunitensi o cinesi, e la superficie
d'attacco offerta da oggetti come termostati o pacemaker. Opinione del docente: il CRA è più importante
di AI Act e NIS2, anche se meno ``di moda''.

\subsection{CRA e NIS2: prodotto contro organizzazione}
\begin{itemize}
  \item \textbf{NIS2} riguarda le reti e le \textbf{aziende}; è una \textbf{direttiva}, attuata in Italia con
    il D.Lgs.~138/2024 (in Italia i soggetti obbligati sono, secondo il docente, circa 20.000).
  \item \textbf{CRA} riguarda il \textbf{prodotto}; è un \textbf{regolamento}, direttamente applicabile in
    tutti gli Stati membri.
\end{itemize}
Obiettivo: tutelare consumatori e imprese che acquistano o usano prodotti o software con una componente
digitale, con requisiti obbligatori per fabbricanti e rivenditori. Problemi affrontati: livello di
sicurezza inadeguato e aggiornamenti insufficienti; impossibilità per gli utenti di capire quali prodotti
siano sicuri. La sicurezza è la prima voce su cui si risparmia per tenere basso il prezzo. Soluzioni:
regole armonizzate per l'immissione sul mercato, requisiti lungo progettazione, sviluppo e manutenzione,
\textbf{dovere di diligenza per l'intero ciclo di vita}. Il considerando~3 spiega perché serve: né il
Cybersecurity Act né la NIS2 coprono direttamente requisiti obbligatori di sicurezza dei prodotti.

\subsection{Ambito di applicazione}
Il CRA si applica ai \textbf{prodotti con elementi digitali messi a disposizione sul mercato} il cui uso
previsto o ragionevolmente prevedibile include una \textbf{connessione dati} diretta o indiretta, logica o
fisica, a un dispositivo o a una rete.
\begin{itemize}
  \item \emph{Prodotto con elementi digitali}: prodotto software o hardware e le relative soluzioni di
    elaborazione dati da remoto, inclusi i componenti immessi sul mercato separatamente.
  \item Connessione \emph{logica} (tramite interfaccia software), \emph{fisica} (cavi, onde radio,
    interfacce elettriche, ottiche, meccaniche), \emph{indiretta} (come parte di un sistema più ampio
    collegabile).
\end{itemize}

\begin{esame}
\textbf{Il software è un prodotto ai fini del CRA.} Errore tipico all'esame: pensare solo a smartwatch e
smart TV ed escludere le applicazioni. Le due parole chiave sono \textbf{mercato} (messa a disposizione
sul mercato, cioè prospettiva commerciale) e \textbf{connessione}.
\end{esame}

\subsection{Esclusioni}
\begin{itemize}
  \item Prodotti già disciplinati da norme settoriali (es.\ regolamento sui \textbf{dispositivi medici}):
    il CRA può essere limitato o escluso se la disciplina settoriale è coerente con il quadro complessivo e
    garantisce un livello di protezione pari o superiore; la Commissione può chiarire con atti delegati.
    Un anello o un orologio smart non è un dispositivo medico; il confine resta discusso (in futuro alcuni
    sistemi di IA generativa potrebbero esservi ricondotti).
  \item \textbf{Pezzi di ricambio} identici, fabbricati secondo le stesse specifiche.
  \item Prodotti sviluppati o modificati esclusivamente per \textbf{sicurezza nazionale o difesa}, o per
    trattare informazioni classificate.
\end{itemize}
Per l'esame basta sapere che esiste un'esclusione generale a favore delle norme settoriali più specifiche.

\subsection{Approccio basato sul rischio e requisiti essenziali}
Il fabbricante deve svolgere un \textbf{cybersecurity risk assessment} del prodotto, che informa tutte le
fasi del ciclo di vita, è documentato e aggiornato durante il periodo di supporto e considera uso previsto e
prevedibile, condizioni d'uso, applicabilità dei requisiti e gestione delle vulnerabilità. Rispetto a un
risk assessment su un processo (es.\ selezione del personale), quello sul dispositivo è
\textbf{profondamente tecnico}: bisogna entrare in ogni processo del dispositivo e nella sua filiera di
fornitori.

\paragraph{Uso prevedibile.} Il fabbricante non può conoscere tutti gli usi futuri; quando emerge un uso non
previsto, una risposta è aggiornare le condizioni d'uso. Il professore cita l'esempio di un assistente di
IA generativa le cui condizioni sono state aggiornate per escludere che si sostituisca a un giudice.

\paragraph{Requisiti per l'immissione sul mercato.} Il prodotto deve rispettare i requisiti essenziali
dell'\textbf{Allegato~I, Parte~I} (se installato, mantenuto e usato correttamente, con gli aggiornamenti
di sicurezza) e i processi del fabbricante quelli dell'\textbf{Allegato~I, Parte~II} (gestione delle
vulnerabilità). Tra i requisiti della Parte~I: nessuna vulnerabilità sfruttabile nota, configurazione
sicura per impostazione predefinita, aggiornamenti di sicurezza, protezione da accessi non autorizzati,
riservatezza e integrità dei dati, \textbf{minimizzazione dei dati}, disponibilità delle funzioni
essenziali, superficie d'attacco ridotta, registrazione dell'attività interna, cancellazione sicura dei dati
da parte dell'utente. Non serve conoscere gli allegati nel dettaglio.

\begin{esame}
\textbf{Spunto per le domande aperte: il rovescio della minimizzazione.} Meno dati significa meno rischio
per l'interessato e per la sicurezza, ma anche meno elementi per il debugging, per l'usabilità e
soprattutto per l'\textbf{analisi forense} dopo un incidente. ``Vince la privacy, muore l'esperto di
cybersecurity'': c'è una tensione da saper argomentare.
\end{esame}

\subsection{Prodotti importanti e critici}
\begin{center}
\small
\begin{tabularx}{\linewidth}{>{\raggedright\arraybackslash}p{3.4cm} X X}
\toprule
& \textbf{Prodotti importanti} & \textbf{Prodotti critici} \\ \midrule
Dove & Allegato~III, classi I e II (es.\ password manager) & Allegato~IV (es.\ dispositivi hardware con
  \emph{security box}), individuati con atti della Commissione \\
Criteri & funzioni critiche per la sicurezza di altri prodotti, reti o servizi, oppure rischio
  significativo di effetti negativi su molti prodotti o sulla salute e sicurezza degli utenti & --- \\
Conseguenza & procedure di valutazione della conformità più rigorose & certificato europeo di
  cybersicurezza con livello di affidabilità almeno \emph{sostanziale} \\ \bottomrule
\end{tabularx}
\end{center}
Il professore ammette che la distinzione pratica non è ancora chiara: per ora l'unico obbligo già operativo
è la \textbf{segnalazione degli incidenti}.

\subsection{Security by design e obblighi lungo la filiera}
\begin{itemize}
  \item \textbf{Fabbricante}: sicurezza integrata in progettazione e sviluppo (come la privacy by design
    del GDPR: la sicurezza è parte della privacy); risk assessment come base di tutto; \textbf{due
    diligence sui componenti di terzi} (supply chain); gestione delle vulnerabilità durante il periodo di
    supporto.
  \item \textbf{Importatore}: verifica che il fabbricante abbia svolto la valutazione di conformità e
    redatto la documentazione tecnica, che ci siano marcatura CE e dichiarazione UE di conformità; azioni
    correttive o ritiro se il prodotto non è conforme; informa fabbricante e autorità di vigilanza del
    mercato delle vulnerabilità scoperte.
  \item \textbf{Distributore}: due diligence, verifica di marcatura CE e documentazione, azioni correttive,
    segnalazione di vulnerabilità e rischi.
\end{itemize}
\paragraph{Il problema pratico.} Il professore cita un grande distributore italiano di hardware e software:
se molti prodotti a basso costo extra UE non rispettano il CRA, dovrà scegliere se toglierli dal catalogo.
Se gli Stati membri applicassero le regole con rigore diverso, i consumatori potrebbero comprare
all'estero ciò che in Italia non si trova più.

\subsection{Obblighi di segnalazione (anticipazione)}
Il fabbricante notifica ogni \textbf{vulnerabilità attivamente sfruttata} al CSIRT tramite la
\textbf{piattaforma unica di segnalazione}, a cui ha accesso anche ENISA: \textbf{early warning entro 24
ore}, \textbf{notifica entro 72 ore}, \textbf{final report entro 14 giorni} dalla disponibilità di una
misura correttiva. Un incidente può far scattare insieme obblighi CRA e di altre norme (24 ore CRA, 72 ore
GDPR\ldots): il tema è ripreso nella lezione del 02/11 (Capitolo~\ref{lez:2026-11-02}) e nella tabella
delle scadenze (Capitolo~\ref{cap:scadenze}).

% =====================================================================
\section{Briefing del laboratorio 3 (Gianluca)}
% =====================================================================
Ruolo: \textbf{responsabile della sicurezza delle informazioni} di un piccolo studio professionale di
consulenza (molto digitale). Il registro contiene già \textbf{18 rischi}: non va inventato nulla.
\begin{itemize}
  \item Per ogni rischio: probabilità e impatto \textbf{inerenti}, separatamente per riservatezza,
    integrità e disponibilità; controlli di trattamento; rischio \textbf{residuo}; decisione finale.
  \item Solo i \textbf{controlli tecnologici} (8.x): niente persone, processi o management.
  \item Matrice \textbf{3$\times$3}; un rischio \textbf{alto} (inerente o residuo) non è accettabile, il
    residuo deve essere medio o basso.
  \item Decisione tra le quattro opzioni; un rischio trasferito resta sotto la responsabilità dell'azienda;
    l'accettazione si usa quando l'inerente è già medio o inferiore.
  \item Dei 34 controlli, 3 non sono applicabili a uno studio di consulenza, 26 sono applicabili e 5 vanno
    valutati caso per caso; ogni esclusione va motivata nelle righe apposite (è ciò che finirebbe nella
    SoA, che però non è oggetto dell'esercizio).
  \item Un foglio di \textbf{coerenza} controlla la completezza formale, non la correttezza nel merito.
  \item Non esiste una risposta unica: \textbf{conta il ragionamento}, non i numeri.
\end{itemize}

\begin{esame}
\textbf{Errore tipico: fare la media} dei punteggi di riservatezza, integrità e disponibilità. Il livello
del rischio è il \textbf{valore massimo} dei tre: un rischio alto solo sulla disponibilità è comunque alto.
\end{esame}

\noindent Svolgimento e metodo completi nel Capitolo~\ref{lab:03}.

% ---------------------------------------------------------------------
\section{Domande aperte e punti da verificare}
\begin{verifica}
\begin{itemize}
  \item \textbf{Threat intelligence nella 27001:2022}: citata come requisito dell'ultima revisione, senza
    numero di controllo nelle slide.
  \item \textbf{Linee guida ENISA per le PMI}: titolo e anno del documento non indicati.
  \item \textbf{Data di applicazione del CRA}: il professore dice ``in vigore da settembre, il 21 se ricordo
    bene''; secondo le slide della lezione del 02/11 gli obblighi di segnalazione si applicano
    dall'\textbf{11 settembre 2026}. Verificare.
  \item \textbf{Circa 20.000 soggetti NIS2 in Italia}: stima a voce del docente.
  \item \textbf{Nuovo decreto legislativo italiano su IA e rapporti di lavoro} (``in vigore da venerdì''),
    citato a proposito dell'uso dell'IA nella selezione del personale: estremi non ricavabili dalla
    trascrizione.
  \item Nomi propri dei casi (Fitbit, Ross Compton, ricercatrice Mandalari) ricostruiti da una trascrizione
    rumorosa: verificare l'ortografia prima di citarli.
  \item Esempi di prodotti importanti e critici: dalle slide (password manager; dispositivi hardware con
    security box); l'elenco completo è negli Allegati III e IV del CRA.
\end{itemize}
\end{verifica}
